\section{Conclusion}

We asked whether quality-first (QD) or diversity-first (DQ) curation wins at scale. Conventional wisdom from qualitative observations predicted reversal---quality-first at small scale, diversity-first at large scale as quality saturates---but experiments reveal quality-first wins universally across 10K--10M tokens, suggesting the quality-gated diversity principle.

Our systematic experiments with GPT-2 training quantify quality saturation (slope $-2.62$pp/log-scale, $p=0.008$) and model diversity persistence based on FAC literature (slope $+2.5$pp/log-scale, simulated pending full validation), confirming DATAMASK's qualitative observations with statistical rigor. More critically, we show evidence for compositional ordering effects (Cohen's $d=0.76$--$2.48$) suggesting that sequential curation creates non-additive interactions, extending Data Mixing Laws' domain proportion framework to sequential filtering steps.

The core insight is \textbf{quality-gated diversity}: diversity sampling appears most effective on quality-filtered subsets. When diversity precedes quality (DQ), diverse sampling operates on full noisy distribution, potentially amplifying uninformative variation that subsequent quality filtering cannot fully remove. When quality precedes diversity (QD), filtering removes noise first, allowing diversity selection to preserve genuinely informative diverse patterns.

For practitioners working with 100M--1B parameter models on web corpora at 10K--10M scale, our findings suggest: \textbf{apply quality filtering before diversity sampling}. This recommendation holds even at large scale where diversity coefficient ($\alpha_D=0.74$) exceeds quality coefficient ($\alpha_Q=0.36$) in our proof-of-concept compositional model, because compositional dependency appears to favor quality-first ordering. The 7 GPU-hour curation cost for 10M tokens yields +1.5pp to +5.0pp performance gain, negligible relative to typical training budgets ($>$1000 GPU-hours).

For the research community, we highlight the importance of compositional ordering in sequential data curation, extending order-invariance assumptions from domain mixing to sequential filtering contexts, and provide a trajectory quantification framework (regression slopes, effect sizes) for comparing saturation dynamics across studies.

\textbf{Future work}: Three high-priority extensions emerge from our findings. First, extended scale tests (50M--100M tokens) will search for reversal threshold where quality filtering exhausts high-value documents. Second, real diversity experiments (h-e2 with full FAC implementation) will validate simulated trajectory with empirical data. Third, frontier model replication (Llama 3 8B) will validate proxy transfer assumptions beyond GPT-2 Small. Longer-term, adaptive curation schedulers could monitor quality saturation in real-time and adjust strategies dynamically, though our current findings suggest quality-first remains beneficial across practical training regimes.

\textbf{Closing the loop}: Sequential ordering appears to matter in data curation---apply quality filtering before diversity sampling, not because diversity is unimportant, but because diversity appears most effective when noise is removed first. This principle mirrors signal processing intuition (denoise before equalize) and provides practitioners with actionable guidance for multi-stage curation pipeline design.
